%! Fundamental Theorem of Algebra and Complex Zeros — Unit 4, Lesson 4.5 — Exit Ticket (Answer Key)
\documentclass[10pt]{article}
\usepackage{algebra2-article}
\usepackage{algebra2-key}   % pulls in -boxes; do NOT also load -boxes
\newcommand{\TallMath}[1]{$\displaystyle #1\rule[-1.4em]{0pt}{3.2em}$}

\begin{document}

\pageheader{Unit 4, Lesson 4.5}{Exit Ticket --- Answer Key}
\namedateperiod

\vspace{0.10in}

No notes. Count your answers against the degree before you turn this in --- and never let an
imaginary zero travel alone.

\begin{enumerate}[leftmargin=1.8em, itemsep=12pt]
  \item \textbf{Solve over the complex numbers.} $x^3 + x^2 + 4x + 4 = 0$

    \vspace{3pt}
    (a) Factor the left side. (Try grouping.) \ans{$x^2(x+1)+4(x+1) = (x+1)(x^2+4)$}

    \vspace{3pt}
    (b) Set each factor equal to zero and solve both: $x = $ \ans{$-1$} \quad and \quad
    $x = $ \ans{$\pm 2i$}

    \vspace{3pt}
    (c) All solutions: \ans{$x = -1,\ 2i,\ -2i$}

    \vspace{3pt}
    (d) \textbf{Number and type:} \ans{$1$} real and \ans{$2$} imaginary. How does the
    degree tell you the list is complete? \ans{Degree $3$ promises exactly $3$ complex zeros, and
    $1+2=3$.}

  \item \textbf{Build one.} Write a polynomial of \emph{least} degree with \emph{real} coefficients
    whose zeros include $-2$ and $4i$.

    \vspace{3pt}
    (a) The full list of zeros: \ans{$-2,\ 4i,\ -4i$} \quad Why? \ans{real coefficients force the
    conjugate $-4i$}

    \vspace{3pt}
    (b) Factored form: \ans{$(x+2)(x^2+16)$}

    \vspace{3pt}
    (c) Standard form: \ans{$x^3 + 2x^2 + 16x + 32$}

  \item \textbf{SOL-style multiple choice.} A polynomial function of degree $4$ has \emph{real}
    coefficients. Which list \emph{cannot} be its complete set of zeros, counted with multiplicity?
    \begin{enumerate}[label=\Alph*., leftmargin=1.9em, itemsep=1pt]
      \item $-1,\; 5,\; 2i,\; -2i$
      \item $3,\; 3,\; -4,\; 0$
      \item $1,\; -1,\; 3i,\; 3i$ \quad \textcolor{keyred}{\textbf{$\leftarrow$ correct}}
      \item $1+i,\; 1-i,\; 2+i,\; 2-i$
    \end{enumerate}
    \begin{itemize}[leftmargin=1.5em, nosep]
      \item \textbf{C} lists $3i$ twice but never lists $-3i$. With real coefficients every imaginary zero needs its conjugate, so this set is impossible.
      \item \textbf{A} pairs $2i$ with $-2i$; \textbf{D} is two complete conjugate pairs; \textbf{B} is all real (a double zero at $3$ is fine --- multiplicity is not the issue).
      \item Each list has four entries, so counting alone does not settle it. The \emph{pairing} is what separates them.
    \end{itemize}
\end{enumerate}

\begin{teachernote}
Three items, three different failure modes, and they sort cleanly.

\vspace{3pt}
\textbf{Item 1} fails if a student writes ``no real solution'' for $x^2+4$ and stops --- the exact
habit this lesson exists to replace --- or drops the $\pm$ and reports only $2i$. Part (d) is the
A2.EI.6d verification and is worth as much as the algebra: a student who solves correctly but cannot
say \emph{why the list is complete} has the procedure without the theorem. Accept any wording of
``degree $3$ means exactly $3$ zeros.''

\vspace{3pt}
\textbf{Item 2} is the conjugate check, and it is pass/fail on one decision. A student who writes
$(x+2)(x-4i)$ has a degree-$2$ answer with an $i$ in its coefficients; ask them to multiply it out.
Watch also for $(x+2)(x^2-16)$ --- the conjugate pair for $\pm 4i$ gives $x^2 + 16$, since
$(x-4i)(x+4i) = x^2 - 16i^2$. That sign is the single most common arithmetic slip of the lesson.

\vspace{3pt}
\textbf{Item 3} tests the theorem as a \emph{filter} rather than a procedure. Distractor B catches
students who think a repeated zero is illegal; distractor D catches students who think ``two
imaginary pairs'' is too many. Students who evaluate or multiply anything here have missed the point
--- the answer is visible from the list alone.

\vspace{3pt}
Sort into four piles: \emph{got it} / \emph{stopped at ``no real solutions''} / \emph{conjugate not
supplied} / \emph{counted wrong}. The conjugate pile is the one to fix immediately, because
Lesson~4.6 has students moving constantly among a graph, a factored form, and a degree, and a missing
conjugate breaks that triangle in all three directions. If the $x^2-16$ sign error is widespread,
sixty seconds re-multiplying $(x-4i)(x+4i)$ on the board at the start of 4.6 is enough.
\end{teachernote}

\end{document}
